\documentclass[11pt]{article}
\usepackage[margin=1in]{geometry}\usepackage{amsmath,amssymb,amsthm,booktabs}
\newtheorem{theorem}{Theorem}\newtheorem{lemma}[theorem]{Lemma}\newtheorem{proposition}[theorem]{Proposition}
\newtheorem{corollary}[theorem]{Corollary}\newtheorem{conjecture}[theorem]{Conjecture}
\theoremstyle{remark}\newtheorem{remark}[theorem]{Remark}
\renewcommand{\Re}{\operatorname{Re}}\newcommand{\osc}{\operatorname{osc}}
\begin{document}
\section*{Idea A: a $p$-uniform route for the relations case ($n=2$, $T=\{1,2\}$)}
\emph{Notation} as in (O1) of \texttt{mechanism\_note.tex}.  $F=\{1,\dots,\frac{p-1}2\}$, $\kappa_l(\delta)=1-\cos(2\pi l\delta/p)=2s_l(\delta)^2$,
$s_l(\delta)=|\sin(\pi l\delta/p)|$, $1\le s\le\frac32$, $r=1/s$, $\tau=3^{s-1}$, $\|a\|=(\sum_la_l^r)^s$, $\Gamma=\Gamma_s(p)$,
$m(a)=\min_{\delta\ne0}(Ka)_\delta$, $\Delta(a)$ its minimisers, and the \emph{efficiency} $e(a)=m(a)/(\Gamma\|a\|)\in[0,1]$.  For $0\le v\le\tau a$ and phases $\theta$,
\[Q(a,v,\theta)=\min_{\delta\ne0}\Big[(Ka)_\delta+c(Kv)_\delta-\tfrac12\osc_wD_\delta\Big],\qquad D_\delta(w)=\sum_lb_l\cos(2\pi lw/p+\psi_l),\ \ b_l=2v_ls_l(\delta),\]
and $(Kv)_\delta=\sum_lb_ls_l(\delta)$.  The relations case is closed iff $Q\le\Gamma\|a\|$ always.  $c_*=1/(2\sqrt3)$ (Theorem A), $c_{12}=1-(2/3)^{s-1}$ (P$_{12}$).

\subsection*{1. Proved: reduction to a weighted Sidon constant}
\paragraph{Definition.}  For $S\subseteq F$ and $\delta\ne0$ let
\[\sigma_\delta(S)=\inf\Big\{\frac{\frac12\osc_w\sum_{l\in S}b_l\cos(2\pi lw/p+\psi_l)}{\sum_{l\in S}b_ls_l(\delta)}\ :\ b\in\mathbb R^S_{\ge0},\ \textstyle\sum b_ls_l(\delta)>0,\ \psi\in\mathbb R^S\Big\}.\]

\begin{theorem}[efficiency--Sidon reduction]\label{es}
Let $a\ge0$, $S=\operatorname{supp}a$, $\delta_0\in\Delta(a)$.  If
\[\sigma_{\delta_0}(S)\ \ge\ c-\frac{1/e(a)-1}{\tau},\]
then $Q(a,v,\theta)\le\Gamma\|a\|$ for every $0\le v\le\tau a$ and all phases.
\end{theorem}
\begin{proof}
Put $\sigma=\sigma_{\delta_0}(S)$; $\operatorname{supp}v\subseteq S$, so $\frac12\osc D_{\delta_0}\ge\sigma(Kv)_{\delta_0}$ and $Q\le m(a)+(c-\sigma)(Kv)_{\delta_0}$.  If $\sigma\ge c$
we are done since $m(a)\le\Gamma\|a\|$.  Otherwise $(Kv)_{\delta_0}\le\tau(Ka)_{\delta_0}=\tau m(a)$ gives $Q\le m(a)(1+(c-\sigma)\tau)\le m(a)/e(a)=\Gamma\|a\|$.
\end{proof}
\begin{corollary}\label{triv}
(a) \emph{(no arithmetic at all)} Every allocation with $e(a)\le1/(1+c\tau)$ is closed.  $c\tau\le\frac12$ for $c_*$ and $c\tau=3^{s-1}-2^{s-1}\le0.318$ for $c_{12}$,
so the thresholds are $e\le\frac23$ resp.\ $e\le0.759$, for every $p$.
(b) \emph{(Parseval)} $\sigma_\delta(S)\ge(\sum_{l\in S}\kappa_l(\delta))^{-1/2}$, so $a$ is closed if $\sum_{S}\kappa_l(\delta_0)\le(c-(1/e(a)-1)/\tau)^{-2}$.
\end{corollary}
\begin{proof}
(b) For mean-zero $P$, $\mathbb EP^2\le(\max P)(-\min P)\le(\frac12\osc P)^2$; $\mathbb EP^2=\frac12\sum b_l^2$ (distinct $l\in F$, $l\not\equiv-l'$); and
$\sum b_ls_l\le(\frac12\sum b_l^2)^{1/2}(2\sum s_l^2)^{1/2}$ with $2\sum s_l^2=\sum\kappa_l$.
\end{proof}
So the whole problem is a lower bound for $\sigma_{\delta_0}(\operatorname{supp}a)$ \emph{only for efficient allocations}, $e(a)>1/(1+c\tau)$, with
the requirement relaxed by the coverage gap.  At $e=1$ (programme maximisers) the requirement is $\sigma\ge c$; Parseval gives it iff
$\sum_{S}\kappa_l(\delta_0)\le1/c^2$.

\subsection*{2. Proved: the tilted-mean bound (any phases)}
\begin{lemma}[tilted mean]\label{tm}
Let $S\subseteq F$, $b\in\mathbb R^S_{\ge0}$, $P(w)=\sum_lb_l\cos(2\pi lw/p+\psi_l)$, $\lambda>0$,
$\Pi(\lambda)=\prod_lI_0(\lambda b_l)$, $Z(\lambda)=\frac1p\sum_{j\in\mathbb Z_p}\exp\big(\lambda\sum_lb_l\cos(2\pi lj/p)\big)$.  Then, for all phases $\psi$,
\[\min(\max P,-\min P)\ \ge\ \frac{2\Pi'(\lambda)-Z'(\lambda)}{Z(\lambda)},\qquad
\min(\max P,-\min P)\ \ge\ \frac{\log\big(2\Pi(\lambda_2)-Z(\lambda_2)\big)-\log Z(\lambda_1)}{\lambda_2-\lambda_1}\ \ (0\le\lambda_1<\lambda_2).\]
Hence $\frac12\osc P$ is at least either right-hand side, and $\sigma_\delta(S)\ge\inf_b\max\{(\frac12\sum b_l^2)^{1/2},\,\sup_\lambda\frac{2\Pi'-Z'}{Z}\}/\sum_lb_ls_l(\delta)$.
\end{lemma}
\begin{proof}
$M(\lambda)=\mathbb E_we^{\lambda P(w)}=\sum_{k\in\mathrm{Rel}(S)}\prod_lI_{k_l}(\lambda b_l)e^{ik\cdot\psi}$, $\mathrm{Rel}(S)=\{k\in\mathbb Z^S:\sum k_ll\equiv0\}$ (Jacobi--Anger and
orthogonality).  The $k=0$ term is $\Pi$; with $\psi=0$ the same sum of moduli is $Z$ (the note's Laplace identity).  Each
$\prod_lI_{|k_l|}(\lambda b_l)$ is a power series in $\lambda$ with non-negative coefficients, so it and its derivative are $\ge0$.  Hence $M\le Z$ and
$M'\ge\Pi'-(Z'-\Pi')$.  $M'/M=\mathbb E_{Q_\lambda}P\le\max P$ for the probability $Q_\lambda\propto e^{\lambda P}$ on $\mathbb Z_p$, which gives the first bound
($\max P\ge0$ anyway).  $\log M$ is convex, so $(\log M(\lambda_2)-\log M(\lambda_1))/(\lambda_2-\lambda_1)\le(\log M)'(\lambda_2)\le\max P$, and
$M(\lambda_2)\ge2\Pi(\lambda_2)-Z(\lambda_2)$, $M(\lambda_1)\le Z(\lambda_1)$.  For $-\min P$ replace $\psi$ by $\psi+\pi$.
\end{proof}
With $x=\lambda b$: $Z/\Pi=1+\rho(x)$, where
\[\rho(x)=\sum_{k\in\mathrm{Rel}(S)\setminus0}\prod_l\frac{I_{|k_l|}(x_l)}{I_0(x_l)}=\frac{\Pr_x(\mathbf K\in\mathrm{Rel}(S)\setminus0)}{\Pr_x(\mathbf K=0)},\qquad
\mathbf K_l\ \text{independent Skellam}(x_l/2,x_l/2),\]
since $e^{-x}I_k(x)$ is the Skellam law.  So $Z(x)e^{-\sum x}$ is the return probability to $0\bmod p$ of the continuous-time walk with jumps $\pm l$ at rate
$x_l/2$, and $\rho(x)=e^{E(x)}/p+\varepsilon(x)$ with $E(x)=\sum_l(x_l-\log I_0x_l)$ and
$\varepsilon(x)=e^{E(x)}\frac1p\sum_{j\ne0}e^{-(Kx)_j}-1$, the excess of the non-zero directions over \emph{independent} phases ($\mathbb Ee^{-\sum x_l\kappa}=e^{-E}$).
Since $I_1/I_0>\log I_0(x)/x$, the first bound is asymptotically $\sum_lb_lA(x_l)$, $A=I_1/I_0$, against $\sum_l\log I_0(x_l)/\lambda$ for the note's
Proposition (tilt): at the programme scale it roughly doubles the certified fraction ($0.386$ vs $0.200$ of $\sum b$ at $K=1.5\ln p$).

\begin{proposition}[conditional, arithmetic-free certificate]\label{cond}
Suppose $(\mathrm{RM}_{\eta,L})$: $\rho(x)\le\eta<1$ for all $x\in\mathbb R^S_{\ge0}$ with $E(x)\le L$.  Then for every $b,\psi$,
\[\min(\max P,-\min P)\ \ge\ \Theta_{\eta,L}(b):=\sup_{\substack{0\le\lambda_1<\lambda_2\\E(\lambda_2b)\le L}}
\frac{\sum_l\log\frac{I_0(\lambda_2b_l)}{I_0(\lambda_1b_l)}-\log\frac{1+\eta}{1-\eta}\mathbf 1_{\lambda_1>0}+\log(1-\eta)\mathbf 1_{\lambda_1=0}}{\lambda_2-\lambda_1},\]
and $\sigma_\delta(S)\ge\theta_{\eta,L}(s(\delta)):=\inf_b\max\{(\frac12\sum b_l^2)^{1/2},\Theta_{\eta,L}(b)\}/\sum b_ls_l(\delta)$, a quantity depending only on the
multiset $\{s_l(\delta)\}_{l\in S}$, $\eta$, $L$ --- not on $p$ or on the arithmetic of $S$.  Consequently (Theorem~\ref{es}) $a$ is closed if $(\mathrm{RM}_{\eta,L})$
holds on $\operatorname{supp}a$ and $\theta_{\eta,L}(s(\delta_0))\ge c-(1/e(a)-1)/\tau$.
\end{proposition}
\begin{proof}
$E$ is coordinatewise increasing ($\frac{d}{dx}(x-\log I_0x)=1-A>0$), so $E(\lambda_1b)\le E(\lambda_2b)\le L$; use $Z\le(1+\eta)\Pi$, $2\Pi-Z\ge(1-\eta)\Pi$ in
Lemma~\ref{tm}, and $M(0)=1$ when $\lambda_1=0$.
\end{proof}
The natural choice is $L=\ln p-\ln(1/\eta_0)$: then $e^{E}/p\le\eta_0$ and $(\mathrm{RM})$ only asks $\varepsilon\le\eta-\eta_0$; with $\varepsilon\le0$ it says that a
non-trivial Skellam vector is a relation of $S$ with probability at most $1/p$, as for independent frequencies.

\begin{proposition}[the $j=0$ barrier]\label{barrier}
For all $S,b,\lambda$: $\frac{2\Pi'-Z'}{Z}\le\min\{2,\,2p\,e^{-E(\lambda b)}\}\sum_lb_lA(\lambda b_l)$, and the secant bound needs $E(\lambda_2b)<\ln(2p)$.  For $b$ uniform on
$K=\kappa\ln p$ frequencies every certificate of Lemma~\ref{tm} is therefore $\le2A(x_\kappa)(1+o(1))\sum b$, $x_\kappa-\log I_0(x_\kappa)=1/\kappa$; this is $\sim\sum b/\kappa$ for large $\kappa$.
\end{proposition}
\begin{proof}
$Z\ge\Pi$, $Z\ge e^{\lambda\sum b}/p$ (the $j=0$ term), $Z'\ge0$, $\Pi'=\Pi\sum bA$.  For uniform $b$ the factor $\min\{2,2pe^{-Kg(x)}\}$, $g(x)=x-\log I_0x$, kills every
$x$ with $Kg(x)\ge\ln p+t$, $t\to\infty$, $t=o(\ln p)$; the secant bound is $\le(\log\Pi)'(\lambda_2)=\sum bA(\lambda_2b)$ by convexity of $\log\Pi$.
\end{proof}
The same $j=0$ argument caps positive-definite product tests $G=\prod_lg_l(2\pi lw/p+\psi_l)$ ($g_l\ge0$, $\hat g_l\ge0$, $\hat g_l(0)=1$):
$\max P\ge(\text{main}-|\text{relation terms}|)/Z_G$ with $Z_G=\frac1p\sum_j\prod_lg_l(2\pi lj/p)\ge\max\{1,\prod_lg_l(0)/p\}$ and main $\le\sum_lb_l\hat g_l(1)\le\sum_lb_l\frac{g_l(0)-1}2$;
for uniform factors this is $\le\frac{e^{1/\kappa}-1}2\sum b\sim\sum b/(2\kappa)$.  So every argument that bounds relation terms by their moduli is capped at order
$(\ln p/K)\sum b$, while the large-deviation value for generic phases is of order $(\ln p/K)^{1/2}\sum b$.  Beating the cap needs cancellation among
relation terms, i.e.\ phase information.

\begin{lemma}[thin frequencies pay for themselves]\label{thin}
Let $a=a'+a''$ with disjoint supports and $a''_l\le\vartheta\|a'\|$ for all $l$, $\vartheta=\big(\Gamma s/(2(1+\tau+c\tau))\big)^{s/(s-1)}$.  If $a'$ is closed
(for all $v'\le\tau a'$, all phases), then $a$ is closed.
\end{lemma}
\begin{proof}
$D_\delta=D'_\delta+D''_\delta$ with $|D''_\delta|\le\sum_lb''_l$, so at the direction $\delta$ that closes $(a',v',\theta)$,
$Q(a,v,\theta)\le\Gamma\|a'\|+\sum_l(a''_l\kappa_l+cv''_l\kappa_l+2v''_ls_l)\le\Gamma\|a'\|+2(1+c\tau+\tau)\sum_la''_l$.  By convexity of $t\mapsto t^s$,
$\Gamma\|a\|\ge\Gamma\|a'\|+\Gamma s\|a'\|^{1-r}\sum_l(a''_l)^r$, and $2(1+c\tau+\tau)a''_l\le\Gamma s\|a'\|^{1-r}(a''_l)^r$ iff $a''_l\le\vartheta\|a'\|$.
\end{proof}
At $s=\frac32$ with $\Gamma\approx\gamma^\ast K^{-1/2}$ and typical core weight $\|a'\|K^{-3/2}$, the threshold is $(0.23\gamma^\ast)^3\approx10^{-3}$ core weights, independent of
$p$; but the number of frequencies above it can still be $\sim10^2K^\ast$, so Theorem~\ref{es} (which ignores the box $v\le\tau a$) must be applied to
a core, not to $\operatorname{supp}a$.

\begin{proposition}[no universal Laplace bound]\label{noUniv}
The identities $\sum_{j\in\mathbb Z_p}\kappa_l(j)=p$ and $\sum_j\kappa_l(j)\kappa_k(j)=p(1+\frac12[l=k])$ ($l,k\in F$) hold for every support.  For $S=\{1,\dots,K\}$ and $x_l\equiv x\in(0,1]$,
$\frac1p\sum_{j\ne0}e^{-(Kx)_j}\ge\big(\frac2{\pi K^{3/2}}-\frac2p\big)e^{-2x}$, while $e^{-E(x)}\le e^{-3Kx/4}$.  Hence no bound for the averaged exponential coverage that uses only
properties every allocation has can give $(\mathrm{RM})$; the hypothesis must use near-optimality ($e(a)>1/(1+c\tau)$).
\end{proposition}
\begin{proof}
$(Kx)_j\le x\sum_{l\le K}2\pi^2l^2j^2/p^2\le2\pi^2xK^3j^2/p^2\le2x$ for $|j|\le p/(\pi K^{3/2})$.  $e^{-x}I_0(x)\le e^{-x+x^2/4}\le e^{-3x/4}$ on $[0,1]$.
\end{proof}
\emph{What survives of ``coverage $\Rightarrow$ Sidon''.}  Min-coverage $C$ bounds every $j\ne0$ term by $e^{-C}$ (the note's form), which needs $C\ge E$, i.e.\ coverage
fraction $\gtrsim0.77$--$0.8$ at the relevant scale; programme maximisers have $0.59$--$0.79$ ($p\le1009$) and the LD model predicts $\to0.46$.  The averaged
form is exactly the Skellam return probability, and that is what has to be controlled.  Transference forces short relations, but their Skellam weight
$\prod_lI_{|k_l|}/I_0\le A(x)^{\|k\|_1}$ is small; what matters is their number, i.e.\ $\varepsilon$.

\subsection*{3. Numerical (files in this directory)}
All at $s=\frac32$ ($b=2$).  Supports: \texttt{supports.json} ($p\le1009$: best of \texttt{relations2/supports\_s1.5.json}, the sibling
\texttt{ideaC} runs, exhaustive search for $p\le29$); $p\ge2003$: local search plus reweighted LP (\texttt{largep.py}, \texttt{refine.py}),
giving near-optimal allocations for several support sizes $K$ per prime.  Efficiencies use the best $\Gamma$ found, which can only overstate $e$
and the $L^2$ ratio; so passes are conservative, failures are not.

\paragraph{N1. Programme maximisers, $5\le p\le1009$} (\texttt{cert.py}, \texttt{cert\_prog.out}, \texttt{relratio.out}).
\begin{center}\small\begin{tabular}{lcccccccccccc}\toprule
$p$ & 29 & 61 & 101 & 151 & 211 & 307 & 401 & 503 & 601 & 701 & 809 & 1009\\
$K^*$ & 4 & 5 & 7 & 6 & 8 & 9 & 8 & 9 & 12 & 10 & 11 & 11\\
$\sum_S\kappa(\delta_0)$ & 3.03 & 3.26 & 4.28 & 3.37 & 4.62 & 5.29 & 4.50 & 4.98 & 7.43 & 5.53 & 5.91 & 5.84\\
$\sigma_P$ & 0.57 & 0.55 & 0.48 & 0.55 & 0.47 & 0.43 & 0.47 & 0.45 & 0.37 & 0.42 & 0.41 & 0.41\\
tilt$/\sum bs$ & 0.42 & 0.56 & 0.50 & 0.60 & 0.44 & 0.41 & 0.53 & 0.49 & 0.40 & 0.44 & 0.44 & 0.49\\
$\max R$ & 0.48 & 0.50 & 0.79 & 0.51 & 0.78 & 0.81 & 0.72 & 0.84 & 0.94 & 0.86 & 0.94 & 0.83\\\bottomrule
\end{tabular}\end{center}
At every one of the 33 primes, $\sigma_{\delta_0}\ge\sigma_P\ge0.367>c_*$, so Theorem~\ref{es} closes the vertex under Theorem A.
The tilted mean at the vertex direction ($b\propto ys$) is comparable to Parseval ($0.30$--$0.60$ for $p\ge23$); its infimum over $b$ is attained near the
Parseval-extremal direction $b\propto s$, so at these sizes it adds $\le0.01$.  Relation mass: $R=\rho(x)/(e^{E(x)}/p)$ at the binding directions and scales
$x_{\max}\in\{\frac12,1,2\}$ is $\le0.97$ for every support ($23\le p\le10007$), median $0.07$--$0.61$; random $K$-sets with the same profile have median
$0.12$--$0.87$ and 90\% quantile up to $1.38$.  So $\varepsilon<0$: programme supports carry fewer weighted relations than independent frequencies.

\paragraph{N2. Spread near-optimal allocations} (\texttt{certlarge.py}, \texttt{cert\_large\_*.out}).  Columns: efficiency $e$, water-filling ratio
$\mathrm{C}_{L2}$ (Theorem A / P$_{12}$; $\le1$ closes), Parseval index and $\sigma_P$, requirement $c-(1/e-1)/\tau$ of Theorem~\ref{es}, and
$\Gamma_{\rm crit}=\min_\delta[(Ka)_\delta+\tau g_\delta(a)]/\|a\|$ (the $L^2$ route closes $a$ iff $\Gamma_s(p)\ge\Gamma_{\rm crit}$).  Best $\Gamma$ found: $0.16425$, $0.14784$,
$0.13999$, $0.12544$ for $p=2003,5003,10007,30011$ (the last two after swap refinement, which raised them by $2.4\%$ and $2.1\%$; the
heuristic still undershoots, so passes are safe and failures are within the uncertainty).  The tilt adds little at these sizes: box-aware tilt
(\texttt{boxtilt.py}) vs water-filling at $p=30011$: $0.973$ vs $0.988$ ($K=23$), $1.006$ vs $1.021$ ($K=26$), $1.032$ vs $1.032$ ($K=30$).  Under P$_{12}$ the tilt alone closes $p=30011$, $K=30$.
(The CMA estimates of $\inf_b$ of the tilt certificate in \texttt{cert\_large\_*.out}, column sigT, are upper estimates and are contradicted by the
box searches; they are not used.)
\begin{center}\small\begin{tabular}{rrcccccccc}\toprule
$p$ & $K$ & $e$ & $\mathrm C_{L2}$(A) & $\mathrm C_{L2}$(P$_{12}$) & $\sum_S\kappa(\delta_0)$ & $\sigma_P$ & need(A) & need(P$_{12}$) & $\Gamma_{\rm crit}$(A)\\\midrule
2003 & 12 & 1.000 & 1.000 & 1.000 & 6.5 & 0.391 & 0.289 & 0.184 & 0.1643\\
2003 & 14 & 0.981 & 0.981 & 0.981 & 7.9 & 0.355 & 0.277 & 0.172 & 0.1611\\
2003 & 16 & 0.969 & 0.969 & 0.969 & 9.7 & 0.321 & 0.270 & 0.165 & 0.1591\\
2003 & 18 & 0.956 & 0.956 & 0.956 & 11.7 & 0.293 & 0.262 & 0.157 & 0.1570\\
2003 & 20 & 0.946 & 0.960 & 0.946 & 13.1 & 0.276 & 0.256 & 0.150 & 0.1576\\
\midrule
5003 & 15 & 0.984 & 0.984 & 0.984 & 8.1 & 0.351 & 0.279 & 0.174 & 0.1455\\
5003 & 17 & 0.990 & 0.990 & 0.990 & 9.9 & 0.317 & 0.283 & 0.178 & 0.1463\\
5003 & 19 & 0.981 & 0.981 & 0.981 & 10.6 & 0.307 & 0.278 & 0.173 & 0.1451\\
5003 & 22 & 0.943 & 0.968 & 0.943 & 13.9 & 0.268 & 0.254 & 0.149 & 0.1430\\
\midrule
10007 & 18 & 0.972 & 0.972 & 0.972 & 10.0 & 0.316 & 0.272 & 0.167 & 0.1360\\
10007 & 20 & 0.963 & 0.963 & 0.963 & 11.6 & 0.294 & 0.267 & 0.161 & 0.1348\\
10007 & 23 & 0.958 & 0.992 & 0.958 & 14.2 & 0.265 & 0.263 & 0.158 & 0.1389\\
10007 & 26 & 0.933 & 0.996 & 0.933 & 16.7 & 0.245 & 0.247 & 0.142 & 0.1395\\
\midrule
30011 & 23 & 0.973 & 0.988 & 0.973 & 13.0 & 0.278 & 0.273 & 0.168 & 0.1240\\
30011 & 26 & 0.975 & 1.021 & 0.975 & 15.6 & 0.254 & 0.274 & 0.169 & 0.1281\\
30011 & 30 & 0.948 & 1.032 & 0.948 & 18.7 & 0.231 & 0.257 & 0.152 & 0.1294\\
\bottomrule\end{tabular}\end{center}


\paragraph{N3. True room} (\texttt{truesig.py}; adversarial search over $(b,\psi)$, calibrated: it finds $0.18$ for the progression $\{1,\dots,26\}$, close
to the Parseval floor $0.14$).  For the spread allocation $p=5003$, $K=22$: $\sigma\le0.69$; for $p=10007$, $K=26$: $\sigma\le0.66$; needed $0.25$--$0.26$, Parseval
gives $0.27$, $0.245$.

\paragraph{N4. Large-deviation model} (\texttt{asym*.py}): $p\to\infty$, $K=\kappa\ln p$, binding-direction law of $\kappa_l(\delta_0)$ = arcsine law tilted to the LD
minimum over $p$ directions, uniform weights, relation excess $\le0$.  The programme optimum of the model is at $\kappa^\ast\approx3.2$, $\gamma^\ast\approx0.46$.
\begin{center}\small\begin{tabular}{lcccccccccccc}\toprule
$\kappa=K/\ln p$ & 2 & 3.2 & 4 & 5 & 6 & 7 & 8 & 9 & 10 & 13 & 17 & 20\\
coverage $\gamma$ & .340 & .464 & .516 & .564 & .600 & .629 & .652 & .671 & .688 & .725 & .759 & .778\\
efficiency $e$ & .928 & 1 & .996 & .973 & .946 & .917 & .889 & .863 & .839 & .776 & .710 & .671\\
need, Thm A & .244 & .289 & .286 & .273 & .255 & .236 & .217 & .197 & .178 & .122 & .053 & .006\\
need, P$_{12}$ & .138 & .184 & .181 & .168 & .150 & .131 & .112 & .092 & .073 & .017 & $<0$ & $<0$\\
tilt, $\inf_b$ & .749 & .388 & .288 & .217 & .175 & .145 & .123 & .107 & .095 & .070 & .052 & .044\\
LD truth ($b\propto s$) & 1.45 & 1.06 & .916 & .797 & .713 & .650 & .600 & .560 & .526 & .452 & .388 & .355\\\bottomrule
\end{tabular}\end{center}
(All entries are fractions of the payment basis $\sum_lb_ls_l(\delta_0)$.)  Under P$_{12}$ the tilt certificate beats the requirement at every $\kappa$
(minimum ratio $1.10$ near $\kappa=8$); under Theorem A it fails for $4\lesssim\kappa\lesssim17$.  The heuristic true value (generic phases) has a
margin $\ge3\times$ everywhere.



\subsection*{4. The exact remaining gap}
\begin{enumerate}
\item[(G1)] \emph{The arithmetic input.}  Prove $(\mathrm{RM}_{\eta,\ln p-L_0})$ for the (cores of) efficient allocations, $e(a)>1/(1+c\tau)$: a non-trivial
Skellam vector on $\operatorname{supp}a$ is a relation mod $p$ with probability $\le(1+\eta')/p$ times the trivial one.  Proposition~\ref{noUniv}
shows that no property shared by all allocations implies it (arithmetic progressions violate it by a factor $p^{\Omega(1)}$); it must come from
near-optimality.  Numerically $R=\rho/(e^E/p)\le0.97$ on every programme support ($23\le p\le10007$), below random sets of the same size.
\item[(G2)] \emph{Theorem A is not enough for any modulus-based argument.}  In the large-deviation model the tilt certificate falls below
$c_*-(1/e-1)/\tau$ for $4\lesssim K/\ln p\lesssim17$ (spread allocations of efficiency $0.71$--$0.99$), and Proposition~\ref{barrier} says no tilt or
positive-definite product test can do better than order $\ln p/K$.  Closing this range with Theorem A needs a bound that uses cancellation among relation
terms (phases), i.e.\ a flat-polynomial theorem for $O(\log p)$-element sets: $\sigma_{\delta_0}(S)\ge c_*-(1/e-1)/\tau$.  Numerically the true $\sigma$ has room
(adversarial search: $\sigma\approx0.66$--$0.69$ against $0.25$--$0.26$ needed), so the inequality is not tight; only the proofs are.
\item[(G3)] \emph{Under P$_{12}$} the model certificate exceeds the requirement for every $K/\ln p$ (minimum ratio $1.10$ at $K/\ln p\approx8$, efficiency
$0.89$).  So P$_{12}$ $+$ (G1) $+$ the model's description of efficient allocations (binding-direction law, uniform weights) gives a $p$-uniform
closure; turning the model into a theorem is the remaining work, and the $10\%$ margin is thin.
\item[(G4)] \emph{Halo.}  Theorem~\ref{es} takes the infimum over all amplitude directions on $\operatorname{supp}a$ and ignores $v\le\tau a$.  Lemma~\ref{thin}
removes frequencies below $\approx10^{-3}$ core weights, $p$-uniformly; between $10^{-3}$ and $1$ core weight a box-aware version (water-filling with the tilt
in place of Parseval) is needed.  It is computable (\texttt{boxtilt.py}) but not yet analysed in the model.
\item[(G5)] \emph{Finite $p$, Theorem A.}  The $L^2$ (water-filling) criterion holds on every allocation tested up to $p=10007$
($\le0.996$ with the best $\Gamma$ found) and fails at $p=30011$ for spread near-optimal allocations ($K=26,30$, $e=0.95$--$0.98$, ratio $1.021$, $1.032$)
unless $\Gamma_s(30011)\ge0.1294$, $3.2\%$ above the best value found ($0.12544$, which was still rising by $\sim1\%$ per refinement round).  The tilt
with exact relation mass helps only marginally (box-aware certificate $0.973$ at $K=23$, $1.006$ at $K=26$, $1.032$ at $K=30$).  So under Theorem A the
relations case is numerically closed up to $p\approx10^4$ and the $L^2$ route is at its edge around $p\approx3\cdot10^4$; beyond, (G2) is the obstruction.
\item[(G6)] \emph{Finite $p$, P$_{12}$.}  $\mathrm C_{L2}\le0.99$ on every allocation tested up to $p=30011$, and at $p=30011$ the tilt \emph{alone}
(no Parseval) already closes the most spread allocation numerically ($K=30$; box search finds no residual gain at three binding directions).  So under P$_{12}$ there is no
visible window between the death of $L^2$ and the takeover by the tilt.
\end{enumerate}

\end{document}
